\section*{الفصل \ref{ch:b2:atomic-orbitals} --- المدارات الذرية: أشكالها وطاقاتها وأحجامها}

\begin{solution}{exo:b2:atomic-orbitals:1}
3p: $n = 3$، $l = 1$؛ \omterm{def:b2:atomic-orbitals:node}{عقدة شعاعية} واحدة ($n - l - 1$) \omterm{def:b2:atomic-orbitals:node}{وعقدة زاوية} واحدة. 4d: $n = 4$،
$l = 2$؛ \omterm{def:b2:atomic-orbitals:node}{عقدة شعاعية} واحدة وعقدتان زاويتان.
\end{solution}

\begin{solution}{exo:b2:atomic-orbitals:2}
$P_{2p} \propto r^4\mathrm e^{-r}$ ($r$ بوحدة $a_0$): $d\ln P/dr = 4/r - 1 = 0$ من أجل $r = 4$:
$4a_0 = \qty{212}{pm}$.
\end{solution}

\begin{solution}{exo:b2:atomic-orbitals:3}
$1 - \mathrm e^{-2}(1 + 2 + 2) = 1 - 5\mathrm e^{-2} = 0.323$.
\end{solution}

\begin{solution}{exo:b2:atomic-orbitals:4}
أربعة فصوص متجهة بين المحورين $x$ و$y$، في المستوي $xy$، متناوبة
الإشارات (موجبة حيث $xy > 0$)؛ والمستويان العقديان $xz$ و$yz$. والعقد الشعاعية:
$n - l - 1 = 0$.
\end{solution}

\begin{solution}{exo:b2:atomic-orbitals:5}
الأكسجين، $(1\mathrm s)^2(2\mathrm s, 2\mathrm p)^6$: $\sigma = 5 \times 0.35 + 2 \times 0.85
= 3.45$، $Z^* = 4.55$؛ ونصف القطر $4/4.55 = 0.88a_0$، أي \qty{47}{pm}.
\end{solution}

\begin{solution}{exo:b2:atomic-orbitals:6}
ينزع التأين الإلكترونات الأقل ارتباطًا. وإلكترونا 4s
(\qty{-14}{eV}) أقل ارتباطًا بكثير من إلكترونات 3d (\qty{-59}{eV}): فيغادر إلكترونا 4s
كلاهما أولًا، ثم إلكترون 3d واحد، فنحصل على $[\ce{Ar}]\,3\mathrm d^5$.
\end{solution}

\begin{solution}{exo:b2:atomic-orbitals:7}
الليثيوم، 2s: $Z^* = 3 - 2 \times 0.85 = 1.30$، $I = 13.6 \times 1.30^2/4 =
\qty{5.75}{eV}$ (المقيسة 5.39، أي أعلى بنسبة 7~\%). والصوديوم، 3s: $Z^* = 11 - 8 \times 0.85 -
2 \times 1.00 = 2.20$، $I = 13.6 \times 2.20^2/9 = \qty{7.3}{eV}$ (المقيسة 5.14): فالقواعد،
وهي نموذج تقريبي، تبالغ في ارتباط الإلكترون الخارجي، أكثر
فأكثر نزولًا في الزمرة.
\end{solution}

\begin{solution}{exo:b2:atomic-orbitals:8}
الصوديوم: $Z^* = 2.20$، ونصف القطر $9/2.20 = 4.1a_0$ (\qty{216}{pm}). والكلور، 3p:
$\sigma = 6 \times 0.35 + 8 \times 0.85 + 2 \times 1.00 = 10.90$، $Z^* = 6.10$، ونصف القطر
$9/6.10 = 1.5a_0$ (\qty{78}{pm}). الطبقة نفسها، لكن الإلكترون الخارجي للكلور
يحس بشحنة تكاد تبلغ ثلاثة أضعاف، فيُجذب نحو الداخل.
\end{solution}

\begin{solution}{exo:b2:atomic-orbitals:9}
لكليهما \omterm{def:b2:atomic-orbitals:radial-angular}{الجزء الزاوي} $1/\sqrt{4\pi}$، فيؤول التكامل إلى
\[
  \int_0^\infty R_{1s}R_{2s}r^2\,dr \propto \int_0^\infty (2 - r)r^2\mathrm e^{-3r/2}\,dr =
  2 \times \frac{2!}{(3/2)^3} - \frac{3!}{(3/2)^4} = \frac{32}{27} - \frac{32}{27} = 0 .
\]
\end{solution}

\begin{solution}{exo:b2:atomic-orbitals:10}
$\int_0^\infty N^2r^2\mathrm e^{-r}r^2\,dr = N^2 \times 4! = 1$: $N = 1/\sqrt{24} =
1/(2\sqrt6)$، كما في الجدول.
\end{solution}

\begin{solution}{exo:b2:atomic-orbitals:11}
$\langle r\rangle = 1.5a_0$، \omterm{def:b2:atomic-orbitals:radial-distribution}{ونصف القطر الأكثر احتمالًا} $a_0$. تصعد $P(r) = 4r^2\mathrm e^{-2r}$
بحدة من الصفر وتنخفض ببطء: فالذيل الطويل عند $r$ الكبيرة يزن في
المتوسط أكثر مما يزن في موضع القيمة العظمى.
\end{solution}

\begin{solution}{exo:b2:atomic-orbitals:12}
$Y_{p_z} \propto z/r = \cos\theta$ ينعدم من أجل $\theta = \pi/2$: المستوي $xy$.
وإشارة $\psi$ هي إشارة $z$، موجبة فوقه وسالبة تحته. وعلى المستوي تكون
\omterm{def:b2:atomic-orbitals:wavefunction}{كثافة الاحتمال} معدومة.
\end{solution}

\begin{solution}{pb:b2:atomic-orbitals:1}
\textbf{1.} $\int|\psi|^2\,dV = \frac{4\pi}{\pi a_0^3}\int_0^\infty r^2\mathrm e^{-2r/a_0}\,dr =
\frac{4}{a_0^3} \times \frac{a_0^3}{4} = 1$.
\textbf{2.} $P(r) = 4\pi r^2|\psi|^2 = (4/a_0^3)\,r^2\mathrm e^{-2r/a_0}$.
\textbf{3.} $dP/dr = 0$: $2/r - 2/a_0 = 0$، $r = a_0$.
\textbf{4.} $\langle r\rangle = (4/a_0^3)\int_0^\infty r^3\mathrm e^{-2r/a_0}\,dr = (4/a_0^3)
\times 6(a_0/2)^4 = 1.5a_0$.
\textbf{5.} التوزيع غير متناظر، وله ذيل طويل عند $r$ الكبيرة.
\textbf{6.} $|\psi|^2$ كثافة لكل وحدة حجم، أكبر ما تكون عند النواة؛ أما $P(r)$
فتحسب حجم القشرة، $4\pi r^2\,dr$، الذي ينعدم عند $r = 0$.
\textbf{7.} \qty{52.9}{pm}.
\textbf{8.} مع $x = r/a_0$: $\int_{x_0}^\infty 4x^2\mathrm e^{-2x}\,dx$؛ وبالتجزئة مرتين،
$= \mathrm e^{-2x_0}(2x_0^2 + 2x_0 + 1)$.
\textbf{9.} $5\mathrm e^{-2} = 0.677$.
\textbf{10.} $13\mathrm e^{-4} = 0.238$.
\textbf{11.} $\mathrm e^{-2x}(1 + 2x + 2x^2) = 0.10$ من أجل $x \approx 2.66$: أي \qty{141}{pm}.
\textbf{12.} لا تنعدم الكثافة أبدًا على أي بعد: فلا حافة. لكن 90~\% من
الاحتمال يقع داخل \qty{141}{pm}: وهذا حجم عملي.
\textbf{13.} 2s: \omterm{def:b2:atomic-orbitals:node}{عقدة شعاعية} واحدة (كرة نصف قطرها $2a_0$)، ولا \omterm{def:b2:atomic-orbitals:node}{عقدة زاوية}. 2p: \omterm{def:b2:atomic-orbitals:node}{عقدة
زاوية} واحدة (مستوٍ)، ولا \omterm{def:b2:atomic-orbitals:node}{عقدة شعاعية}.
\textbf{14.} عند $r = 2a_0$.
\textbf{15.} $4a_0$ (التمرين 2).
\textbf{16.} $d\ln P/dr = 2/r - 2/(2 - r) - 1 = 0$ تعطي $r^2 - 6r + 4 = 0$، $r = 3 \pm \sqrt5$:
$0.76a_0$ (قيمة عظمى داخلية صغيرة) و$5.24a_0$ (القيمة العظمى الرئيسية).
\textbf{17.} 2s، عبر قيمته العظمى الداخلية. ففي ذرة متعددة الإلكترونات ينفذ إلكترون 2s
عبر طبقة 1s، فيُحجب أقل من إلكترون 2p ويقع أدنى في
الطاقة.
\textbf{18.} تُقسم أنصاف الأقطار كلها على $Z = 6$: 2p عند $0.67a_0$ (\qty{35}{pm})؛ والقيمتان العظميان للمدار 2s
عند $0.13a_0$ و$0.87a_0$.
\textbf{19.} C 3.25، N 3.90، O 4.55.
\textbf{20.} $\varepsilon = -13.6\,Z^{*2}/4$: C \qty{-35.9}{eV}، N \qty{-51.7}{eV}، O
\qty{-70.4}{eV}.
\textbf{21.} كما في الفصل: \qty{11.5}{eV} (المقيسة 11.26).
\textbf{22.} N: الذرة $5 \times (-13.6 \times 3.90^2/4) = \qty{-258.6}{eV}$؛ \ce{N+}، $Z^* =
4.25$، $4 \times (-13.6 \times 4.25^2/4) = \qty{-245.7}{eV}$؛ $I = \qty{12.9}{eV}$.
\textbf{23.} O: الذرة $6 \times (-13.6 \times 4.55^2/4) = \qty{-422.3}{eV}$؛ \ce{O+}، $Z^* =
4.90$، $5 \times (-13.6 \times 4.90^2/4) = \qty{-408.2}{eV}$؛ $I = \qty{14.2}{eV}$. النموذج
11.5 و12.9 و14.2؛ والمقيسة 11.26 و14.53 و13.62.
\textbf{24.} الطاقة المقيسة للأكسجين أدنى من طاقة النيتروجين: ففي
الأكسجين يجب أن يتقاسم إلكترون 2p الرابع مدارًا مع إلكترون آخر، 
وتنافرهما يجعل نزعه أسهل، بينما تكوّن إلكترونات p الثلاثة المنفردة في النيتروجين
طبقة فرعية نصف ممتلئة مستقرة. أما \omterm{def:b2:atomic-orbitals:slater}{قواعد سلاتر} فتعامل كل إلكترونات المجموعة
بالطريقة نفسها ولا تعرف شيئًا عن الاقتران.
\textbf{25.} يوجد إلكترون 1s في الهيدروجين أبعد من $2a_0$ باحتمال
$\boldsymbol{13\mathrm e^{-4} \approx 0.238}$.
\end{solution}
